\documentclass[11pt]{article}
\newcommand{\doctitle}{Deriving the Biomarker--Treatment Correlation
  Decay from First Principles}
\newcommand{\shorttitle}{BM-BR Correlation Decay}
\newcommand{\docdate}{2026-03-20}
% pmsimstats-ng standard document preamble
% Usage: % pmsimstats-ng standard document preamble
% Usage: % pmsimstats-ng standard document preamble
% Usage: \input{pmsimstats-preamble} after \documentclass[11pt]{article}
%
% Documents must define before \input:
%   \newcommand{\doctitle}{...}       Full document title
%   \newcommand{\shorttitle}{...}     Short title for header
%   \newcommand{\docdate}{YYYY-MM-DD} Document date

\usepackage[margin=1in]{geometry}
\usepackage{amsmath,amssymb}
\usepackage[T1]{fontenc}
\usepackage{lmodern}
\usepackage{microtype}
\usepackage{graphicx}
\graphicspath{{figures/}}
\usepackage{booktabs}
\usepackage{longtable}
\usepackage{array}
\usepackage{hyperref}
\usepackage{xcolor}
\usepackage{fancyhdr}
\usepackage{parskip}
\usepackage{caption}
\usepackage{enumitem}
\usepackage{verbatim}
\usepackage{listings}

\lstset{
  language=R,
  basicstyle=\small\ttfamily,
  keywordstyle=\color{blue!70!black},
  commentstyle=\color{green!50!black},
  stringstyle=\color{red!60!black},
  breaklines=true,
  frame=single,
  framerule=0.5pt,
  rulecolor=\color{gray!40},
  backgroundcolor=\color{gray!5},
  xleftmargin=1em,
  framexleftmargin=1em,
  columns=flexible
}

\hypersetup{
  colorlinks=true,
  linkcolor=blue!60!black,
  citecolor=green!50!black,
  urlcolor=blue!60!black
}

\pagestyle{fancy}
\fancyhf{}
\lhead{\footnotesize pmsimstats-ng Technical Report}
\rhead{\footnotesize \shorttitle}
\rfoot{\footnotesize \thepage}
\renewcommand{\headrulewidth}{0.4pt}

\title{\doctitle}
\author{pmsimstats team}
\date{\docdate}
 after \documentclass[11pt]{article}
%
% Documents must define before \input:
%   \newcommand{\doctitle}{...}       Full document title
%   \newcommand{\shorttitle}{...}     Short title for header
%   \newcommand{\docdate}{YYYY-MM-DD} Document date

\usepackage[margin=1in]{geometry}
\usepackage{amsmath,amssymb}
\usepackage[T1]{fontenc}
\usepackage{lmodern}
\usepackage{microtype}
\usepackage{graphicx}
\graphicspath{{figures/}}
\usepackage{booktabs}
\usepackage{longtable}
\usepackage{array}
\usepackage{hyperref}
\usepackage{xcolor}
\usepackage{fancyhdr}
\usepackage{parskip}
\usepackage{caption}
\usepackage{enumitem}
\usepackage{verbatim}
\usepackage{listings}

\lstset{
  language=R,
  basicstyle=\small\ttfamily,
  keywordstyle=\color{blue!70!black},
  commentstyle=\color{green!50!black},
  stringstyle=\color{red!60!black},
  breaklines=true,
  frame=single,
  framerule=0.5pt,
  rulecolor=\color{gray!40},
  backgroundcolor=\color{gray!5},
  xleftmargin=1em,
  framexleftmargin=1em,
  columns=flexible
}

\hypersetup{
  colorlinks=true,
  linkcolor=blue!60!black,
  citecolor=green!50!black,
  urlcolor=blue!60!black
}

\pagestyle{fancy}
\fancyhf{}
\lhead{\footnotesize pmsimstats-ng Technical Report}
\rhead{\footnotesize \shorttitle}
\rfoot{\footnotesize \thepage}
\renewcommand{\headrulewidth}{0.4pt}

\title{\doctitle}
\author{pmsimstats team}
\date{\docdate}
 after \documentclass[11pt]{article}
%
% Documents must define before \input:
%   \newcommand{\doctitle}{...}       Full document title
%   \newcommand{\shorttitle}{...}     Short title for header
%   \newcommand{\docdate}{YYYY-MM-DD} Document date

\usepackage[margin=1in]{geometry}
\usepackage{amsmath,amssymb}
\usepackage[T1]{fontenc}
\usepackage{lmodern}
\usepackage{microtype}
\usepackage{graphicx}
\graphicspath{{figures/}}
\usepackage{booktabs}
\usepackage{longtable}
\usepackage{array}
\usepackage{hyperref}
\usepackage{xcolor}
\usepackage{fancyhdr}
\usepackage{parskip}
\usepackage{caption}
\usepackage{enumitem}
\usepackage{verbatim}
\usepackage{listings}

\lstset{
  language=R,
  basicstyle=\small\ttfamily,
  keywordstyle=\color{blue!70!black},
  commentstyle=\color{green!50!black},
  stringstyle=\color{red!60!black},
  breaklines=true,
  frame=single,
  framerule=0.5pt,
  rulecolor=\color{gray!40},
  backgroundcolor=\color{gray!5},
  xleftmargin=1em,
  framexleftmargin=1em,
  columns=flexible
}

\hypersetup{
  colorlinks=true,
  linkcolor=blue!60!black,
  citecolor=green!50!black,
  urlcolor=blue!60!black
}

\pagestyle{fancy}
\fancyhf{}
\lhead{\footnotesize pmsimstats-ng Technical Report}
\rhead{\footnotesize \shorttitle}
\rfoot{\footnotesize \thepage}
\renewcommand{\headrulewidth}{0.4pt}

\title{\doctitle}
\author{pmsimstats team}
\date{\docdate}


\begin{document}
\maketitle
\thispagestyle{fancy}

\begin{abstract}
The \texttt{pmsimstats} data-generating process draws simulated
trial data from a multivariate normal distribution whose
correlation structure encodes a biomarker--treatment interaction.
How the biomarker--biologic response (BM--BR) correlation should
behave during off-drug periods with carryover has been a central
methodological question, with three different implementations
across the repository's history. This paper derives the
correlation decay rule from first principles by tracing what
happens to individual participant drug effects during
pharmacokinetic clearance. The derivation shows that the
correlation should decay at the same exponential rate as the
mean drug effect, preserving the coefficient of variation of
the biomarker-specific component. This result validates the
core insight of the Ron Thomas proportional scaling while
providing a principled parameterization via
$\lambda_{cor} = \ln 2 / t_{1/2}$.
\end{abstract}

\tableofcontents
\newpage

% =================================================================
\section{The Problem}
% =================================================================

The \texttt{pmsimstats} simulation generates participant outcomes
by drawing from a multivariate normal (MVN) distribution. At each
timepoint $t$, the biologic response (BR) factor has three
properties specified in the MVN:

\begin{enumerate}
\item \textbf{Mean} $\mu_{\text{BR},t}$: the expected drug
  effect, computed from the modified Gompertz function and
  carryover decay.
\item \textbf{Standard deviation} $\sigma_{\text{BR}} = 8$:
  the between-person variability in drug response, constant
  across all timepoints.
\item \textbf{Correlation with biomarker}
  $\rho_t = \text{Cor}(\text{BM}, \text{BR}_t)$: determines
  how much the biomarker predicts individual variation in
  drug response.
\end{enumerate}

During on-drug periods, $\rho_t = c_{bm}$ (the specified
biomarker correlation). During off-drug periods with carryover,
the mean decays but the variance remains fixed. The question is:
what should $\rho_t$ be?

Three approaches have been implemented:

\begin{description}
\item[Step function (original, commit \texttt{42ac030}):]
  $\rho_t = c_{bm}$ if $\mu_{\text{BR},t} \neq 0$, else 0.
  Full correlation at any nonzero residual, regardless of
  magnitude.

\item[Proportional scaling (Ron Thomas, commit
  \texttt{8609f12}):] $\rho_t = (\mu_{\text{BR},t} /
  \mu_{\text{BR},t-1}) \cdot c_{bm}$. Correlation decays
  proportionally with the mean ratio between adjacent
  timepoints.

\item[Independent exponential (current HEAD):]
  $\rho_t = c_{bm} \cdot e^{-\lambda_{cor} \cdot t_{sd}}$.
  Correlation decays exponentially with time since
  discontinuation, with $\lambda_{cor}$ as a free parameter.
\end{description}

This paper asks: is there a principled way to determine
$\rho_t$ from the structure of the problem, rather than
treating it as an arbitrary modeling choice?

% =================================================================
\section{The Data-Generating Process}
% =================================================================

\subsection{The MVN Framework}

Each participant $i$ receives a draw from the MVN. Focusing on
the biomarker and BR factor at a single timepoint $t$:

\[
\begin{pmatrix} \text{BM}_i \\ \text{BR}_{i,t} \end{pmatrix}
\sim \mathcal{N}\!\left(
\begin{pmatrix} \mu_{\text{BM}} \\ \mu_{\text{BR},t}
\end{pmatrix},
\begin{pmatrix}
\sigma_{\text{BM}}^2 &
  \rho_t \sigma_{\text{BM}} \sigma_{\text{BR}} \\
\rho_t \sigma_{\text{BM}} \sigma_{\text{BR}} &
  \sigma_{\text{BR}}^2
\end{pmatrix}
\right)
\]

The conditional distribution of BR given a participant's
biomarker value is:

\begin{equation}
\text{BR}_{i,t} \mid \text{BM}_i = bm \;\sim\;
\mathcal{N}\!\left(
  \mu_{\text{BR},t} +
  \rho_t \frac{\sigma_{\text{BR}}}{\sigma_{\text{BM}}}
  (bm - \mu_{\text{BM}}),\;
  \sigma_{\text{BR}}^2 (1 - \rho_t^2)
\right)
\label{eq:conditional}
\end{equation}

The conditional mean has two components:

\begin{itemize}
\item $\mu_{\text{BR},t}$: the population-average drug effect
  (same for everyone)
\item $\rho_t (\sigma_{\text{BR}} / \sigma_{\text{BM}})
  (bm - \mu_{\text{BM}})$: the biomarker-specific deviation
  (varies by participant)
\end{itemize}

\subsection{The Gompertz Mean Trajectory}

The BR mean at on-drug timepoints follows the modified Gompertz:

\[
\mu_{\text{BR},t}^{\text{on}} =
  f_G(\text{tod}_t;\; 10.99, 5, 0.42)
\]

At off-drug timepoints, the raw Gompertz evaluates to zero
($\text{tod} = 0$), and carryover provides the mean:

\[
\mu_{\text{BR},t}^{\text{off}} =
  \mu_{\text{BR},t-1} \cdot
  \left(\tfrac{1}{2}\right)^{t_{sd} / t_{1/2}}
= \mu_{\text{BR},t-1} \cdot \alpha
\]

where $\alpha = (1/2)^{t_{sd}/t_{1/2}}$ is the decay factor.

\subsection{Parameter Values}

\begin{table}[h]
\centering
\caption{Key parameter values from the pilot RCT data}
\begin{tabular}{llrl}
\toprule
Parameter & Symbol & Value & Source \\
\midrule
Biomarker mean & $\mu_{\text{BM}}$ & 124.3 mmHg &
  Pilot systolic BP \\
Biomarker SD & $\sigma_{\text{BM}}$ & 15.4 mmHg &
  Pilot systolic BP \\
BR SD & $\sigma_{\text{BR}}$ & 8.0 CAPS points &
  Active$-$placebo variability \\
Biomarker correlation & $c_{bm}$ & 0--0.5 &
  Simulation parameter \\
\bottomrule
\end{tabular}
\end{table}

% =================================================================
\section{Derivation from Pharmacokinetic Principles}
% =================================================================

\subsection{The Individual-Level Argument}

Consider two participants during an on-drug period when the
population BR mean is $\mu$:

\begin{itemize}
\item Participant A has biomarker 1 SD above the population
  mean ($bm_A = \mu_{\text{BM}} + \sigma_{\text{BM}}$).
\item Participant B has biomarker 1 SD below the population
  mean ($bm_B = \mu_{\text{BM}} - \sigma_{\text{BM}}$).
\end{itemize}

From Equation~\ref{eq:conditional}, their expected BR values
are:

\begin{align}
E[\text{BR}_A] &= \mu + \rho \cdot \sigma_{\text{BR}} = \mu + c_{bm} \cdot 8 \label{eq:brA} \\
E[\text{BR}_B] &= \mu - \rho \cdot \sigma_{\text{BR}} = \mu - c_{bm} \cdot 8 \label{eq:brB}
\end{align}

The biomarker-specific spread is:

\begin{equation}
\Delta_{\text{on}} = E[\text{BR}_A] - E[\text{BR}_B]
= 2 \cdot c_{bm} \cdot \sigma_{\text{BR}}
\label{eq:spread_on}
\end{equation}

\subsection{What Happens at Drug Discontinuation?}

When the drug is discontinued, the pharmacologic effect decays.
If carryover is pharmacokinetic (drug molecules clear at a
rate determined by the drug's half-life, independent of the
patient's biomarker value), then \emph{each participant's
drug effect decays by the same fraction} $\alpha$:

\begin{align}
E[\text{BR}_A^{\text{off}}] &= \alpha \cdot E[\text{BR}_A]
  = \alpha \cdot (\mu + c_{bm} \cdot 8) \\
E[\text{BR}_B^{\text{off}}] &= \alpha \cdot E[\text{BR}_B]
  = \alpha \cdot (\mu - c_{bm} \cdot 8)
\end{align}

The population mean at the off-drug timepoint is:

\[
\mu_{\text{off}} = \alpha \cdot \mu
\]

which matches the carryover formula in the DGP. The
biomarker-specific spread is:

\begin{equation}
\Delta_{\text{off}} = E[\text{BR}_A^{\text{off}}] -
  E[\text{BR}_B^{\text{off}}]
= \alpha \cdot 2 \cdot c_{bm} \cdot \sigma_{\text{BR}}
= \alpha \cdot \Delta_{\text{on}}
\label{eq:spread_off}
\end{equation}

The spread decays by the same factor $\alpha$ as the mean.

\subsection{Implied Correlation}

From Equation~\ref{eq:conditional}, the biomarker-specific
spread at the off-drug timepoint is also:

\[
\Delta_{\text{off}} = 2 \cdot \rho_{\text{off}} \cdot
  \sigma_{\text{BR}}
\]

Setting this equal to Equation~\ref{eq:spread_off}:

\[
2 \cdot \rho_{\text{off}} \cdot \sigma_{\text{BR}}
= \alpha \cdot 2 \cdot c_{bm} \cdot \sigma_{\text{BR}}
\]

\begin{equation}
\boxed{\rho_{\text{off}} = \alpha \cdot c_{bm}
= \left(\tfrac{1}{2}\right)^{t_{sd}/t_{1/2}} \cdot c_{bm}}
\label{eq:result}
\end{equation}

The BM--BR correlation at off-drug timepoints should decay by
the same exponential factor as the BR mean. This is the central
result.

\subsection{Equivalence to Exponential Parameterization}

The decay factor $\alpha$ can be written as:

\[
\alpha = \left(\tfrac{1}{2}\right)^{t_{sd}/t_{1/2}}
= e^{-(\ln 2 / t_{1/2}) \cdot t_{sd}}
= e^{-\lambda_{drug} \cdot t_{sd}}
\]

where $\lambda_{drug} = \ln 2 / t_{1/2}$ is the drug's
elimination rate constant. Substituting into
Equation~\ref{eq:result}:

\[
\rho_{\text{off}} = c_{bm} \cdot
  e^{-\lambda_{drug} \cdot t_{sd}}
\]

This matches the current HEAD parameterization
$\rho_t = c_{bm} \cdot e^{-\lambda_{cor} \cdot t_{sd}}$
when:

\begin{equation}
\boxed{\lambda_{cor} = \lambda_{drug} = \frac{\ln 2}{t_{1/2}}}
\label{eq:lambda}
\end{equation}

The correlation decay rate should equal the drug elimination
rate constant.

% =================================================================
\section{Verification of Internal Consistency}
% =================================================================

\subsection{Coefficient of Variation}

The coefficient of variation (CV) of the biomarker-specific
component measures how informative the biomarker is relative
to the magnitude of the drug effect:

\[
\text{CV} = \frac{\Delta}{\mu} =
\frac{2 \cdot \rho \cdot \sigma_{\text{BR}}}{\mu}
\]

On drug: $\text{CV}_{\text{on}} = 2 \cdot c_{bm} \cdot
\sigma_{\text{BR}} / \mu$

Off drug with proportional scaling:
$\text{CV}_{\text{off}} = 2 \cdot (\alpha \cdot c_{bm})
\cdot \sigma_{\text{BR}} / (\alpha \cdot \mu)
= 2 \cdot c_{bm} \cdot \sigma_{\text{BR}} / \mu
= \text{CV}_{\text{on}}$

The CV is preserved. The biomarker is equally informative
\emph{relative to the remaining effect}, which is the
pharmacokinetically consistent expectation.

\subsection{Numerical Example}

N-of-1 Path A, $c_{bm} = 0.5$, $t_{1/2} = 1.0$ week:

\begin{table}[h]
\centering
\caption{Conditional expectations under proportional
  correlation decay}
\begin{tabular}{lrrrrrr}
\toprule
& \multicolumn{3}{c}{On drug (BD2, week 10)} &
  \multicolumn{3}{c}{Off drug (BD3, week 11)} \\
\cmidrule(lr){2-4} \cmidrule(lr){5-7}
& Mean & $E[\text{BR}|+1$SD$]$ & $E[\text{BR}|-1$SD$]$ &
  Mean & $E[\text{BR}|+1$SD$]$ & $E[\text{BR}|-1$SD$]$ \\
\midrule
$\mu_{\text{BR}}$ & 10.19 & & & 5.09 & & \\
$\rho$ & 0.50 & & & 0.25 & & \\
$E[\text{BR}|\text{BM}]$ & & 14.19 & 6.19 & & 7.09 & 3.09 \\
Spread & & \multicolumn{2}{c}{8.0} & &
  \multicolumn{2}{c}{4.0} \\
CV & & \multicolumn{2}{c}{0.785} & &
  \multicolumn{2}{c}{0.786} \\
\bottomrule
\end{tabular}
\end{table}

The spread halves (8.0 to 4.0) as the mean halves (10.19 to
5.09). The CV is preserved at 0.79. No participant has a
negative expected BR. The conditional expectations are
clinically plausible.

\subsection{Comparison with Alternative Rules}

\begin{table}[h]
\centering
\caption{Off-drug conditional expectations at BD3
  ($\mu_{\text{BR}} = 5.09$, $c_{bm} = 0.5$) under
  three correlation rules}
\begin{tabular}{lrrrl}
\toprule
Rule & $\rho$ & $E[\text{BR}|+1$SD$]$ &
  $E[\text{BR}|-1$SD$]$ & Issue \\
\midrule
Step function ($\lambda_{cor} = 0$) & 0 & 5.09 & 5.09 &
  No biomarker signal \\
Proportional ($\lambda_{cor} = \lambda_{drug}$) & 0.25 &
  7.09 & 3.09 & \textbf{Consistent} \\
Slow decay ($\lambda_{cor} = 1$) & 0.18 & 6.57 & 3.61 &
  Under-predicts spread \\
Full correlation ($\rho = c_{bm}$) & 0.50 & 9.09 & 1.09 &
  Over-predicts spread \\
\bottomrule
\end{tabular}
\end{table}

Only the proportional scaling produces a spread that is
consistent with the pharmacokinetic decay of individual
effects. The step function ($\lambda_{cor} = 0$) eliminates
the biomarker signal entirely. Full correlation
($\rho = c_{bm}$) maintains the on-drug spread despite the
mean halving, implying the biomarker is twice as informative
per unit of drug effect during carryover as during active
treatment.

% =================================================================
\section{Implications for the Publication's Findings}
% =================================================================

\subsection{The Original Power Drop Was an Artifact}

The original code used the step-function rule at short
carryover half-lives ($t_{1/2} = 0.1, 0.2$ weeks). At these
values, the carryover residual is negligible ($< 0.1\%$ of
peak), but nonzero. The step function assigns full $c_{bm}$
correlation to these timepoints, producing large
biomarker-specific variation around a near-zero mean. This
injects spurious noise into the interaction estimate.

Under proportional scaling, the correlation at these
timepoints would be $\alpha \cdot c_{bm} \approx 0.001
\cdot c_{bm} \approx 0$, which is effectively the same as
no correlation. The power drop disappears because the
artifact disappears.

\subsection{At Realistic Half-Lives, Power Should Drop
            Moderately}

At clinically realistic half-lives ($t_{1/2} = 0.5$--$1.0$
weeks), the carryover residual is substantial (25--50\% of
peak) and the proportional correlation is meaningful
($\rho \approx 0.13$--$0.25 \cdot c_{bm}$). The off-drug
timepoints carry partial biomarker interaction information
that partially compensates for the carryover confounding.

The net effect on power depends on the balance between:

\begin{enumerate}
\item \textbf{Confounding (reduces power):} The carryover
  residual in the mean structure blurs the on-drug/off-drug
  contrast.
\item \textbf{Preserved signal (increases power):} The
  proportional correlation maintains biomarker-specific
  information during carryover periods.
\end{enumerate}

Our 200-replicate simulation with $\lambda_{cor} = 1$
(close to $\lambda_{drug} = 0.693$ at $t_{1/2} = 1.0$)
showed flat power across carryover values, suggesting these
two effects approximately cancel at the tested parameters.

\subsection{The Correct $\lambda_{cor}$ Value}

From Equation~\ref{eq:lambda}, the correct value of
$\lambda_{cor}$ is not a fixed constant but depends on the
carryover half-life being simulated:

\begin{table}[h]
\centering
\caption{$\lambda_{cor}$ values for different carryover
  half-lives}
\begin{tabular}{rr}
\toprule
$t_{1/2}$ (weeks) & $\lambda_{cor} = \ln 2 / t_{1/2}$ \\
\midrule
0.1 & 6.93 \\
0.2 & 3.47 \\
0.5 & 1.39 \\
1.0 & 0.693 \\
2.0 & 0.347 \\
\bottomrule
\end{tabular}
\end{table}

A fixed $\lambda_{cor}$ (such as the current default of 1.0)
is only correct for $t_{1/2} \approx 0.693$ weeks
($\approx 4.8$ days). For other half-lives, it either
over-decays the correlation (short $t_{1/2}$) or
under-decays it (long $t_{1/2}$). The implementation
should compute $\lambda_{cor}$ from the carryover half-life
rather than accepting it as a separate parameter.

% =================================================================
\section{Relationship to the Three Historical Implementations}
% =================================================================

\subsection{Step Function (Original)}

The step function is the limiting case
$\lambda_{cor} \to \infty$ for timepoints where $\mu = 0$
and $\lambda_{cor} = 0$ for timepoints where $\mu > 0$.
It is a discontinuous approximation that fails when the
mean is small but nonzero (the carryover regime). It
produces clinically implausible conditional expectations
(negative BR for low-biomarker participants) and
artifact-driven power loss.

\subsection{Ron Thomas Proportional Scaling}

The proportional scaling computes $\rho_t = (\mu_t /
\mu_{t-1}) \cdot c_{bm}$. Since $\mu_t / \mu_{t-1} =
\alpha$ (the single-step decay factor), this is equivalent
to $\rho_t = \alpha \cdot c_{bm}$, which is exactly
Equation~\ref{eq:result}. The Ron Thomas implementation
is the pharmacokinetically correct solution.

However, the implementation had three problems:

\begin{enumerate}
\item It used a ratio of adjacent means rather than the
  decay factor directly, requiring a $p > 1$ guard that
  inadvertently zeroed the correlation at timepoint 1.
\item It was bundled with a scale factor of 2 on the
  carryover decay that changed the meaning of
  $t_{1/2}$.
\item It was undocumented and unjustified in the commit.
\end{enumerate}

\subsection{Independent Exponential (Current HEAD)}

The current parameterization $\rho_t = c_{bm} \cdot
e^{-\lambda_{cor} \cdot t_{sd}}$ is the correct functional
form. Setting $\lambda_{cor} = \ln 2 / t_{1/2}$ recovers
the pharmacokinetically derived result. The advantage over
the Ron Thomas implementation is that it avoids the ratio
formula, handles timepoint 1 correctly, and is explicit
about the decay mechanism.

% =================================================================
\section{Recommended Implementation}
% =================================================================

The \texttt{generateData()} function should compute
$\lambda_{cor}$ automatically from the carryover half-life
rather than accepting it as a separate parameter:

\begin{verbatim}
# In generateData():
if (modelparam$carryover_t1half > 0) {
  lambda_cor <- log(2) / modelparam$carryover_t1half
} else {
  lambda_cor <- 0
}
\end{verbatim}

This ensures that the correlation decay is always consistent
with the mean decay, regardless of the half-life being
simulated. The separate $\lambda_{cor}$ parameter can be
retained as an override for sensitivity analysis (testing
scenarios where the biomarker's predictive power decays
faster or slower than the drug effect), but the default
should be the pharmacokinetically derived value.

The BM--BR correlation assignment then becomes:

\begin{verbatim}
for (p in 1:nP) {
  n1 <- paste(trialdesign$timeptname[p], "br", sep = ".")
  if (d[p]$onDrug) {
    correlations[n1, 'bm'] <- modelparam$c.bm
    correlations['bm', n1] <- modelparam$c.bm
  } else if (d[p]$tsd > 0 && lambda_cor > 0) {
    decay <- exp(-lambda_cor * d$tsd[p])
    correlations[n1, 'bm'] <- modelparam$c.bm * decay
    correlations['bm', n1] <- modelparam$c.bm * decay
  }
}
\end{verbatim}

% =================================================================
\section{The Remaining Variance Problem}
% =================================================================

The derivation in Section~3 assumes that the BR standard
deviation $\sigma_{\text{BR}} = 8$ is constant, which is
the current implementation. This creates a situation where
the conditional variance
$\sigma_{\text{BR}}^2 (1 - \rho_t^2)$ is nearly unchanged
during off-drug periods (since $\rho_t$ is small), meaning
participants still have large random fluctuations in BR
around a decaying mean.

This is a deeper structural limitation of the MVN
parameterization. The SD of 8 conflates two sources of
variation:

\begin{enumerate}
\item \textbf{Trait variability:} Between-person differences
  in latent drug responsiveness, which exist before the drug
  is administered and persist regardless of drug status. This
  component should remain constant.

\item \textbf{State-dependent variability:} Variation in
  the acute drug effect, which should scale with the
  magnitude of the effect. This component should decay during
  carryover.
\end{enumerate}

The current framework cannot separate these two sources.
The proportional correlation scaling addresses the
biomarker-specific component correctly (the spread decays
with the mean), but the total variance remains inflated
during off-drug periods. This means the conditional
distribution has the correct center (proportional spread)
but excessive width (fixed SD). The practical consequence
is that the simulation adds more noise during off-drug
periods than is clinically realistic, which may
conservatively underestimate power.

Addressing this limitation would require a time-varying
$\sigma_{\text{BR},t}$ or a decomposition of the BR
variance into trait and state components, either of which
represents a substantial restructuring of the DGP. For the
current application (comparing relative power across trial
designs), the fixed-variance assumption is acceptable
because it affects all designs similarly.

% =================================================================
\section{Conclusions}
% =================================================================

\begin{enumerate}
\item The BM--BR correlation during off-drug carryover
  periods should decay at the same exponential rate as the
  BR mean: $\rho_t = c_{bm} \cdot (1/2)^{t_{sd}/t_{1/2}}$.
  This follows from the pharmacokinetic principle that drug
  clearance affects all participants' effects by the same
  fraction, preserving the coefficient of variation of the
  biomarker-specific component.

\item This result is equivalent to setting
  $\lambda_{cor} = \ln 2 / t_{1/2}$ in the current
  exponential parameterization, linking the correlation
  decay to the drug elimination rate constant.

\item The Ron Thomas proportional scaling (commit
  \texttt{8609f12}) implemented the correct principle but
  with an unnecessarily complex formula and three
  co-packaged changes (scale factor, timepoint-1 skip)
  that introduced new problems.

\item The recommended implementation computes
  $\lambda_{cor}$ automatically from
  \texttt{carryover\_t1half}, making the correlation decay
  consistent with the mean decay by construction, while
  retaining $\lambda_{cor}$ as an optional override for
  sensitivity analysis.

\item The power drop reported in Hendrickson et al.\ (2020)
  at $t_{1/2} = 0.1$ and $0.2$ weeks was an artifact of the
  step-function correlation rule, not a genuine consequence
  of carryover at those half-lives. At clinically realistic
  half-lives ($t_{1/2} = 0.5$--$1.0$ weeks), the
  proportional correlation decay partially compensates for
  carryover confounding, producing moderate rather than
  catastrophic power loss.

\item The fixed BR variance ($\sigma_{\text{BR}} = 8$) is a
  known limitation that conservatively inflates noise during
  off-drug periods. It does not affect the validity of the
  correlation decay derivation but may underestimate power
  relative to a model with time-varying variance.
\end{enumerate}

\vfill
\noindent\rule{\textwidth}{0.4pt}
{\footnotesize
Rendered on 2026-03-25 at 18:49 PDT.\\
Source: \textasciitilde/prj/alz/10-pmsimstats-ng/pmsimstats-ng/docs/08-biomarker-correlation-decay.tex}

\end{document}
